\subsection{Lysogorskiy et al.\ (2021): PACE applied to Cu and Si}
\label{sec:appendix-example-lysogorskiy2021ace}

\paragraph{Q1 -- Bibliographic identification.}
Lysogorskiy, van der Oord, Bochkarev, Menon, Rinaldi, Hammerschmidt,
Mrovec, Thompson, Cs\'anyi, Ortner, and Drautz introduce PACE, a
performant C++ implementation of the Atomic Cluster Expansion
(\cite{drautz2019ace}) released as a LAMMPS pair-style. Production
parameterisations are presented for fcc copper and diamond silicon.
Published in \emph{npj Computational Materials} 7, 97 (2021);
doi:10.1038/s41524-021-00559-9.
\lstinputlisting[language=turtle]{artifacts/kg/listings/lysogorskiy2021ace/q01-bibliographic.ttl}

\paragraph{Q2 -- Material system.}
Two single-element systems are encoded as separate
$\MaterialSystemC$ instances: copper (fcc, with bcc/dhcp/hcp
polymorphs and surfaces/defects in the reference set) and silicon
(diamond cubic, with reference data from the GAP-Si general-purpose
database of Bart\'ok et al.\ 2018).
\lstinputlisting[language=turtle]{artifacts/kg/listings/lysogorskiy2021ace/q02-material.ttl}

\paragraph{Q3 -- Reference calculation method.}
Both reference datasets come from DFT calculations.
\lstinputlisting[language=turtle]{artifacts/kg/listings/lysogorskiy2021ace/q03-reference-method-type.ttl}

\paragraph{Q4 -- Reference settings.}
Cu reference data are produced with FHI-aims at the PBE level
(small clusters, bulk phases, surfaces and slabs, configurations
with displaced or missing atoms), partly assembled with the
\texttt{pyiron} workflow framework. Si reference data are the
published GAP-Si general-purpose fitting database from CASTEP
(PW91 inherited from Bart\'ok et al.\ 2018).
\lstinputlisting[language=turtle]{artifacts/kg/listings/lysogorskiy2021ace/q04-reference-settings.ttl}

\paragraph{Q5 -- Training dataset.}
Cu: in-house FHI-aims set covering clusters, bulk polymorphs,
surfaces, slabs, and defects (configuration count not enumerated in
the main text). Si: the published Bart\'ok-2018 GAP-Si database of
2{,}475 configurations covering crystal phases, surfaces, vacancies,
interstitials, and liquid Si. Cu covers energies+forces; Si
additionally covers virials.
\lstinputlisting[language=turtle]{artifacts/kg/listings/lysogorskiy2021ace/q05-dataset.ttl}

\paragraph{Q6 -- Sampling strategies.}
Cu: configurational coverage of clusters, polymorphs, surfaces, and
defective structures. Si: reuse of the published GAP-Si database's
sampling strategy.
\lstinputlisting[language=turtle]{artifacts/kg/listings/lysogorskiy2021ace/q06-sampling.ttl}

\paragraph{Q7 -- MLIP method, components, implementation.}
The method is the Atomic Cluster Expansion of Drautz~2019: per-atom
property $\varphi_i$ expanded over body-ordered multi-atom basis
functions, with truncation at body order $\nu_{\max}$ and total
energy as a sum over per-atom embeddings. The Cu fit uses a
nonlinear two-density Finnis--Sinclair embedding; the Si fit uses
a linear embedding at $\nu_{\max}=4$. Implementation: PACE
(C++ LAMMPS \texttt{pair\_style}).
\lstinputlisting[language=turtle]{artifacts/kg/listings/lysogorskiy2021ace/q07-method.ttl}

\paragraph{Q8 -- Hyperparameter settings.}
Reported settings: Cu cutoff radius $r_c = 7.4$\,\AA; Cu basis size
$2{,}072$ parameters (756 expansion coefficients per atomic density
$\times$ 2 densities + 560 radial-function parameters); Si linear
basis size $6{,}827$ functions at body order $\nu_{\max}=4$.
\lstinputlisting[language=turtle]{artifacts/kg/listings/lysogorskiy2021ace/q08-hyperparameter-settings.ttl}

\paragraph{Q9 -- MLIP run and trained model.}
Two training runs are encoded: one for PACE-Cu (nonlinear
Finnis--Sinclair embedding) and one for PACE-Si (linear ACE,
$\nu_{\max}=4$). Both produce a single $\TrainedModelC$ with a
recorded per-atom inference time
($\dlRole{inferenceTimePerAtom}=0.00032$\,s/atom for Cu,
$0.0008$\,s/atom for Si).
\lstinputlisting[language=turtle]{artifacts/kg/listings/lysogorskiy2021ace/q09-run-and-model.ttl}

\paragraph{Q10 -- Benchmark results.}
PACE-Cu attains a final fit error of 2.9\,meV/atom on structures
within 1\,eV of the ground state (after ground-state fine-tuning),
with force errors near 15\,meV/\AA{} on higher-energy structures.
PACE-Si matches GAP-Si accuracy on the Bart\'ok-2018 database with
energy MAE 1.81\,meV/atom and force MAE 82\,meV/\AA. Three
$\BenchmarkResultC$ instances cover these.
\lstinputlisting[language=turtle]{artifacts/kg/listings/lysogorskiy2021ace/q10-benchmark-results.ttl}

\paragraph{Q11 -- Computational resources.}
The headline claim is computational performance: PACE shifts the
accuracy/cost Pareto front for ML interatomic potentials. We capture
the per-atom inference times on the trained models in Q9
(0.32\,ms/atom for Cu, 0.80\,ms/atom for Si on a single force
call). Training duration, GPU/CPU hours, training hardware, and
peak memory are not reported.
\lstinputlisting[language=turtle]{artifacts/kg/listings/lysogorskiy2021ace/q11-resources.ttl}

\paragraph{Q12 -- Gaps.}
The protocol's "one paper $\equiv$ one method $\times$ one material"
framing is mildly stretched: the paper trains \emph{two} ACE
parameterisations (Cu and Si) and we encode both within one
\texttt{lysogorskiy2021ace.ttl} via two $\MaterialSystemC$, two
$\dlConcept{MLIPRun}$, and two $\TrainedModelC$ instances rather
than splitting into sibling files. Other unreported items: explicit
DFT settings on the Cu side beyond engine and functional (energy
cutoff, k-mesh, pseudopotential / NAO basis); explicit
$\dlRole{numConfigurations}$ for the Cu reference set; the Pareto
benchmark is reported as a derived figure rather than an
$\AccuracyMetricC$ across all five models compared (ACE, EAM,
SNAP, GTINV, GAP); training-cost metadata.
